\section*{TD 3 --- Régression et régularisation}
\addcontentsline{toc}{section}{TD 3 --- Régression et régularisation}

\subsection*{Exercice 1 --- Ajuster une droite à la main \hfill \normalfont\footnotesize (12 min)}
Quatre points : $(1,3)$, $(2,5)$, $(3,6)$, $(4,8)$.
\begin{enumerate}
  \item Calculer $\bar x$ et $\bar y$.
  \item Calculer la pente $\hat\beta_1=\dfrac{\sum (x_i-\bar x)(y_i-\bar y)}{\sum (x_i-\bar x)^2}$.
  \item En déduire l'ordonnée à l'origine, puis la prédiction en $x=5$.
  \item Calculer le $R^2$.
\end{enumerate}

\subsection*{Exercice 2 --- Interpréter un coefficient \hfill \normalfont\footnotesize (6 min)}
Un modèle estime $\hat y = 15{,}2 + 0{,}42\,x$ où $x$ est la surface en m$^2$
et $y$ le loyer en milliers de FCFA.
\begin{enumerate}
  \item Que signifie le coefficient $0{,}42$ ?
  \item Que signifie l'ordonnée $15{,}2$ ? A-t-elle un sens ici ?
  \item Prédire le loyer pour 60 m$^2$.
\end{enumerate}

\subsection*{Exercice 3 --- Ridge ou Lasso ? \hfill \normalfont\footnotesize (8 min)}
On dispose de 40 variables, dont beaucoup sont redondantes.
\begin{enumerate}
  \item Quelle pénalité privilégier si l'on veut \textbf{sélectionner} des
        variables ? Pourquoi ?
  \item Laquelle privilégier si l'on veut seulement \textbf{stabiliser} les
        coefficients sans en annuler ?
  \item Comment choisit-on $\lambda$ ?
\end{enumerate}

\subsection*{Exercice 4 --- RMSE contre MAE \hfill \normalfont\footnotesize (8 min)}
Deux modèles, sur le même jeu de test :
\begin{center}
\small
\begin{tabular}{@{}lcc@{}}
\toprule
& \textbf{RMSE} & \textbf{MAE} \\
\midrule
Modèle A & 12,0 & 11,4 \\
Modèle B & 15,2 & 8,1 \\
\bottomrule
\end{tabular}
\end{center}
\begin{enumerate}
  \item Lequel a les plus grosses erreurs ponctuelles ?
  \item Lequel choisir si une grosse erreur coûte très cher ?
  \item Lequel choisir si l'on veut l'erreur typique la plus faible ?
\end{enumerate}

\subsection*{Exercice 5 --- Diagnostic de résidus \hfill \normalfont\footnotesize (8 min)}
On observe un nuage de résidus en forme de U en fonction de la prédiction.
\begin{enumerate}
  \item Quel défaut du modèle cela révèle-t-il ?
  \item Citer deux remèdes possibles.
  \item Que signifierait au contraire un entonnoir (variance croissante) ?
\end{enumerate}


\subsection*{Exercice 6 --- Interpréter un $R^2$ \hfill \normalfont\footnotesize (6 min)}
Deux modèles obtiennent $R^2=0{,}98$ et $R^2=0{,}45$ sur le même problème.
\begin{enumerate}
  \item Que signifie chaque valeur en termes de variance expliquée ?
  \item Le $R^2$ élevé garantit-il de bonnes prévisions ? Pourquoi ?
  \item Que ne dit pas le $R^2$ ?
\end{enumerate}

\subsection*{Exercice 7 --- Multicolinéarité \hfill \normalfont\footnotesize (7 min)}
Deux variables sont corrélées à $0{,}95$ et leurs coefficients changent de signe
d'un échantillon à l'autre.
\begin{enumerate}
  \item Nommer le phénomène.
  \item Pourquoi les coefficients deviennent-ils instables ?
  \item Citer deux remèdes.
\end{enumerate}

\subsection*{Exercice 8 --- Une descente de gradient à la main \hfill \normalfont\footnotesize (10 min)}
On minimise $J(w)=(w-3)^2+1$, avec $w_0=0$ et $\eta=0{,}1$.
\begin{enumerate}
  \item Calculer le gradient en $w_0$.
  \item En déduire $w_1$, puis $w_2$ et $w_3$.
  \item Vers quelle valeur la suite converge-t-elle ?
\end{enumerate}

\subsection*{Exercice 9 --- MAE ou MSE ? \hfill \normalfont\footnotesize (6 min)}
Les données de prix contiennent quelques valeurs extrêmes.
\begin{enumerate}
  \item Quelle perte est la plus robuste ? Pourquoi ?
  \item Quelle perte pénalise davantage une grosse erreur ponctuelle ?
  \item Laquelle choisir si le coût est proportionnel à l'erreur ?
\end{enumerate}

\newpage
